\BLOCK{ if ts_extra_packages }
\BLOCK{- set ns = namespace(has_tabularx=False, has_marginnote=False) }
\BLOCK{ for name, options in ts_extra_packages }
\BLOCK{- if name == "tabularx" -}
\BLOCK{- set ns.has_tabularx = True -}
\BLOCK{- endif }
\BLOCK{- if name == "marginnote" -}
\BLOCK{- set ns.has_marginnote = True -}
\BLOCK{- endif }
\usepackage\BLOCK{ if options }[\VAR{options}]\BLOCK{ endif }{\VAR{name}}
\BLOCK{ endfor }
\BLOCK{- if ns.has_tabularx }
\newcolumntype{Y}{>{\centering\arraybackslash}X}
\newcolumntype{Z}{>{\raggedleft\arraybackslash}X}
\BLOCK{- endif }
\BLOCK{- if ns.has_marginnote }
% Default margin-note styling: footnote-size so notes fit narrow margins
% without overflowing. Override with ``\renewcommand*{\marginfont}{…}`` in a
% custom preamble snippet if you need a different look.
\renewcommand*{\marginfont}{\footnotesize}
% Clamp \marginparwidth so margin notes never overflow the page edge, even
% when the user's geometry reserves less horizontal space than LaTeX's
% default \marginparwidth (~3.9cm). We compute the space actually available
% on each side (recto / verso) at \AtBeginDocument time — once geometry has
% finalised \textwidth / \oddsidemargin / \evensidemargin — and keep the
% smaller of the two so the note fits on both sides in twoside documents.
\newlength{\tsmarginparavail}
\newlength{\tsmarginparalt}
\makeatletter
\AtBeginDocument{%
  \setlength{\tsmarginparavail}{\dimexpr\paperwidth-\textwidth-1in-\oddsidemargin-\marginparsep\relax}%
  \ifdim\tsmarginparavail<\z@\setlength{\tsmarginparavail}{\z@}\fi
  \if@twoside
    \setlength{\tsmarginparalt}{\dimexpr1in+\evensidemargin-\marginparsep\relax}%
    \ifdim\tsmarginparalt<\z@\setlength{\tsmarginparalt}{\z@}\fi
    \ifdim\tsmarginparalt<\tsmarginparavail
      \setlength{\tsmarginparavail}{\tsmarginparalt}%
    \fi
  \fi
  \ifdim\tsmarginparavail<\marginparwidth
    \setlength{\marginparwidth}{\tsmarginparavail}%
  \fi
}
\makeatother
\BLOCK{- endif }
\BLOCK{ endif }
